% CausalLine: selective recovery after prompt injection in multi-agent LLM systems
%
% Compiles with any standard TeX distribution (TeX Live, MiKTeX, Overleaf).
% Only widely available packages are used; no custom class file is required.
%
%   pdflatex causalline && pdflatex causalline
%
\documentclass[10pt,twocolumn]{article}

\usepackage[T1]{fontenc}
\usepackage[utf8]{inputenc}
\usepackage[margin=0.75in]{geometry}
\usepackage{amsmath}
\usepackage{amssymb}
\usepackage{booktabs}
\usepackage{graphicx}
\usepackage{xcolor}
\usepackage{array}
\usepackage[hidelinks]{hyperref}

\newcommand{\sys}{\textsc{CausalLine}}
\newcommand{\bzero}{\textsf{B0}}
\newcommand{\bone}{\textsf{B1}}
\newcommand{\btwo}{\textsf{B2}}

\title{\bf Exposure Is Not Influence:\\
Selective Recovery After Prompt Injection\\
in Multi-Agent LLM Systems}

\author{%
Anonymous Submission
}

\date{}

\begin{document}
\maketitle

\begin{abstract}
Work on prompt injection in multi-agent LLM systems overwhelmingly addresses
\emph{detection}. We study what happens next. Once an agent is known to be
compromised, current practice discards that agent's output and everything
downstream of it. In a 56-agent workflow with a broadcast hub, that is
$85.5\%$ of the execution trace for a single poisoned record.

We observe that this over-discards because it confuses two different
relations: a source being \emph{present} in an agent's context (exposure) and
a source \emph{changing} what the agent produced (influence). \sys{} records
provenance during execution, establishes influence per (source, event) pair,
propagates contamination only along influence and derivation edges, replays
only the affected events, and---critically---\emph{verifies} the recovered run
before accepting it, escalating to agent restart and then full restart when
verification fails.

Across 15 runs of a real 56-agent, three-provider system (50 local
Llama-3-8B, 2 Gemini, 4 Nemotron) under three attack channels and three
contamination regimes, \sys{} preserved strictly more correct work than the
topology-closure baseline in 13 runs, tied in 2, and lost in 0, while
recording \textbf{zero unsafe preservations, zero contamination escapes, and
recall $1.0$}. On a deterministic scripted matrix (96 cells $\times$ 30
repetitions, four detectors including a blind control), \sys{} recorded zero
unsafe preservations in all 24 of its cells.

We are equally explicit about what this costs. \sys{} is \emph{more expensive
in tokens than every baseline we compare against}, including full restart:
$1.6\times$--$2.4\times$ the cost of simply re-running the workflow. We
therefore do not claim it is an optimization. We give the exact condition
under which it is the right choice---when completed work cannot be
re-executed, because it has side effects, external cost, or a non-deterministic
outcome that a restart may not reproduce---and we recommend restart when it
can. We further report two negative results that shaped the work: the
falsification of our own ex-ante economic gate, and a provenance defect whose
correction removed an apparent $+17.2$ point performance win that was, in
fact, a safety failure.
\end{abstract}

\section{Introduction}

A multi-agent LLM system is a pipeline of agents that read from tools, memory,
the web, and each other. Any of those channels can carry an injected
instruction~\cite{greshake2023,perez2022,liu2024}. A large and growing
literature studies how to \emph{notice} this.

Far less attention is paid to the step after. Suppose detection works: an
operator learns at 14:03 that the Researcher agent consumed a poisoned page.
What should be thrown away?

The answer in practice is: that agent's output, and everything downstream of
it. This is safe and it is enormously wasteful. It is wasteful because
\emph{reachability is not causation}. An agent that received a poisoned value
in its context may never have used it. A hub that broadcasts a roster of
twelve records to twelve specialists exposes all twelve to a poisoned record,
while each specialist uses exactly one.

\paragraph{The gap, measured.}
In our 56-agent testbed the distinction is not subtle. Table~\ref{tab:gap}
shows the fraction of the execution trace that the topology closure discards
against the fraction that was actually influenced.

\begin{table}[t]
\centering
\small
\begin{tabular}{lrrr}
\toprule
regime & structural closure & true influence & gap \\
\midrule
small  & $7.6\%$  & $6.2\%$  & $1.4$ pts \\
medium & $60.7\%$ & $54.5\%$ & $6.2$ pts \\
large  & $85.5\%$ & $72.4\%$ & $13.1$ pts \\
\bottomrule
\end{tabular}
\caption{Exposure versus influence on a 145-event, 56-agent trace. The
closure is what reachability-based recovery discards; true influence is what
a canary token shows was actually affected.}
\label{tab:gap}
\end{table}

\paragraph{Contributions.}
\begin{enumerate}
\itemsep2pt
\item \textbf{A recovery method, not a detector.} \sys{} takes a detector's
verdict as input and produces a minimal, verified recovery
(Section~\ref{sec:design}).
\item \textbf{A soundness discipline that survives adversarial reporting.}
Anything not \emph{established} clean is treated as contaminated; a source an
agent falsely claims not to have used still contaminates unless a
counterfactual says otherwise (Section~\ref{sec:attribution}).
\item \textbf{Verification with escalation.} A recovery is never shipped
unverified. Failure escalates selective $\rightarrow$ agent restart
$\rightarrow$ full restart (Section~\ref{sec:verify}).
\item \textbf{Evaluation on a real heterogeneous system}: 56 agents, six
stages, three inference providers, three attack channels, three contamination
regimes, 15 runs (Section~\ref{sec:eval}).
\item \textbf{An honest economic characterisation}, including the finding
that \sys{} costs more tokens than full restart, and the exact condition under
which it is nonetheless correct to use (Section~\ref{sec:econ}).
\item \textbf{Two negative results} we consider as important as the positive
ones (Section~\ref{sec:negative}).
\end{enumerate}

\section{Problem Setting}

\subsection{Threat model}

An adversary controls the content of at least one \emph{source}: a web page a
tool returns, a value stored in shared memory, or a message from one agent to
another. The adversary cannot modify our code, the trace log, or the recovery
machinery. We assume an external detector eventually identifies some set of
malicious sources; \sys{} consumes that verdict and does not produce it. In
the evaluation we use an oracle detector in the real-system experiments---so
that what we measure is recovery quality and not detection quality---and four
detectors of varying strength, including a blind control that flags nothing,
in the scripted experiments.

\subsection{Vocabulary}

We use these terms precisely throughout.

\begin{description}
\itemsep1pt
\item[event] one recorded operation: a model call, a tool call, a tool
response, a memory read or write, an inter-agent message.
\item[source] an incoming unit of information, with an identifier.
\item[exposure] a source was present in an agent's context.
\item[influence] a source demonstrably changed an event's output.
\item[contaminated] influenced, directly or transitively, by a malicious
source.
\item[unsafe preservation] recovery kept an event that was in fact
contaminated. This is the dangerous error.
\end{description}

\subsection{Baselines}

\begin{description}
\itemsep1pt
\item[\bzero{} full restart] discard everything, re-run. Preserves $0\%$ by
construction. Always safe.
\item[\bone{} agent-level taint] discard every event produced by an agent that
touched a flagged source, and everything downstream.
\item[\btwo{} topology closure] discard everything reachable in the call graph
from where a flagged source entered. This is the strongest baseline and the
one we compare against throughout.
\end{description}

\section{Design}
\label{sec:design}

\sys{} has five stages. Only the first runs during the workflow; the rest run
after a detector fires.

\subsection{Execution-time provenance}

Every event is logged with its agent, its parents, the identifiers of every
source in its context (\emph{exposure}, whether or not it was used), and a
content-addressed reference to what it produced. Prompts and outputs are
stored so that a counterfactual can be re-issued later.

Two record kinds are written without any model in the loop, because the
answer is a property of the code path rather than an estimate: a
\emph{structural} record for operations our own code computes, and a
\emph{carrier} record for an event whose output is a copy of an earlier
event's, which inherits that event's verdicts rather than asserting its own.

\paragraph{Ingestion edges.}
An event that \emph{retrieves} a source stores that source's text in its own
output. The source does not exist when the event is logged, so it is not in
that event's exposures, so no verdict is ever written for the pair. We record
an explicit ingestion edge from each materialised source back to the event
that materialised it. Section~\ref{sec:negative} describes what happened
before we did.

\subsection{Influence attribution}
\label{sec:attribution}

For each (source, event) pair the estimator produces one of: \emph{influenced},
\emph{clean}, or \emph{unchecked}. Two mechanisms are combined.

\textbf{Self-report} asks the agent which of its inputs changed its output. It
is cheap and it is not trusted: a self-report \emph{positive} is accepted
(claiming to have used something is not a claim that buys the agent anything),
while a self-report \emph{negative} is only a hypothesis.

\textbf{Counterfactual replay} re-issues the event's stored prompt with the
source redacted and compares outputs. This is evidence, and it is what clears
a pair.

The asymmetry is the safety property. A \emph{tainted} verdict grows the
recovery region and is therefore sound to accept on weak evidence; a
\emph{clean} verdict shrinks it and requires strong evidence. Anything
unchecked is treated as contaminated.

Leave-one-out is only valid when redaction actually removes the information.
We check \emph{removability}: if the same fact survives redaction through
another route, the check is void and the pair stays contaminated.

\subsection{Contamination propagation}

From the flagged sources, contamination walks exactly three relations:
influence (source $\rightarrow$ event), derivation (event $\rightarrow$ the
source wrapping its output for a downstream agent), and ingestion (source
$\rightarrow$ the event that retrieved it).

It never walks parent edges. Parent edges are execution history, not
causation; walking them transitively reproduces \btwo{} exactly. A method that
silently reproduces its own baseline has no result left to report, and the
failure is invisible because the numbers still come out.

\subsection{Recovery planning and selective replay}

The planner computes the invalidation set and the nearest trusted checkpoint
per agent. Replay re-executes only invalidated events; every other event is
spliced from its logged output, byte for byte. Two invariants are asserted:
every spliced event's output matches its original log entry, and no model call
is issued for an event outside the invalidation set. Flagged sources are
redacted from re-issued prompts, and a redaction that fails aborts the
recovery rather than producing a run that reports a removal that did not
happen.

\subsection{Verification and escalation}
\label{sec:verify}

A recovered run is checked before it is accepted: the task-level outcome must
hold, no live memory entry may still point at an invalidated write, and no
flagged source may remain in a re-issued prompt. On failure, recovery escalates:

\[
\text{selective} \rightarrow \text{agent restart} \rightarrow \text{full restart}
\]

Restart is therefore never the primary action and always the guaranteed
fallback. This matters for the claims we can make: \sys{} cannot be worse than
restart in \emph{outcome}, only in \emph{cost}.

\section{Implementation}

\sys{} is roughly 9{,}000 lines of Python with no dependencies outside the
standard library for the core, and is instrumented to log every token it
spends against the purpose it spent it on. The evaluation harness supports two
modes that we never mix, because their ground truth is different in kind:

\textbf{Scripted mode} drives the pipeline with a deterministic client whose
usage rule we wrote, so ground truth is known \emph{by construction} for every
(source, event) pair. This is the only mode in which a per-pair accuracy
number is meaningful.

\textbf{Real-LLM mode} runs on hosted and local models. Ground truth is
\emph{observed}: the injected payload instructs the agent to emit a specific
wrong value, and that value is a canary. If it appears in an output, that
output was influenced.

\section{Evaluation}
\label{sec:eval}

\subsection{The 56-agent testbed}

Prior evaluations in this line of work use pipelines of four or five agents on
one model. That cannot support a claim about a deployed system. We therefore
built a 56-agent workflow, summarised in Table~\ref{tab:topology}, spanning
three inference providers.

\begin{table}[t]
\centering
\small
\begin{tabular}{llr}
\toprule
stage & role & agents \\
\midrule
A & acquisition (each fetches its own record) & 12 \\
B & normalisation (memory + a block of A)     & 10 \\
C & \textbf{hub} (fan-in over B, broadcasts)  & 1 \\
C & specialists (hub roster + one B block)    & 12 \\
D & verifiers (a block of specialists)        & 10 \\
E & reviewers (a block of D)                  & 8 \\
F & synthesis, audit, executor                & 3 \\
\midrule
  & \textbf{total} & \textbf{56} \\
\bottomrule
\end{tabular}
\caption{The testbed workflow: six stages, 145 events. Providers: 50 agents
on local Llama-3-8B, 2 on Gemini-3.8-Flash (the two fan-in points), 4 on
Nemotron.}
\label{tab:topology}
\end{table}

The hub is the reason the topology exists. It reproduces every record it
receives and every specialist reads that roster, so a poisoned source reaching
the hub is \emph{exposed} to the entire downstream trace while each specialist
\emph{uses} exactly one record. If fine-grained attribution cannot beat a
reachability closure here, it cannot anywhere.

Provider placement is not decoration: the web-channel contamination path runs
local $\rightarrow$ Gemini $\rightarrow$ Nemotron $\rightarrow$ local, so
every hop is a different model with its own tokenizer and its own
instruction-following behaviour.

\subsection{Attack channels}

Three channels, chosen because they propagate differently:

\begin{enumerate}
\itemsep1pt
\item \textbf{web/tool}: a note beside a fetched record. Arrives through a
tool call, so it has an origin event.
\item \textbf{poisoned memory}: a policy entry a normaliser reads. \emph{A
memory entry has no call-graph edge back to whoever wrote it}, so this is the
case where a cross-agent flow exists without a topology edge.
\item \textbf{inter-agent message}: planted into a verifier's context with no
derivation link---it is not the output of any event in the run.
\end{enumerate}

\subsection{Contamination regimes}

Three regimes, defined by where the attack enters, with the resulting
contaminated fraction \emph{measured} from each trace rather than configured.

\subsection{Protocol}

Five seeds $\times$ three regimes $=15$ runs, all completed. The architecture,
provider assignment, payloads, regime definitions, baselines, \sys{}
configuration, metrics and detector were fixed before the first run and not
touched afterwards. All four methods are scored on the \emph{same} trace under
the \emph{same} attack in each run, so the comparison is paired.

\section{Results}

\subsection{Preserved work}

\begin{table*}[t]
\centering
\small
\begin{tabular}{lrrrrrrrrr}
\toprule
regime & $f_{\text{true}}$ & \bzero{} & \bone{} & \btwo{} & \sys{} &
$\Delta$ vs \btwo{} & sd & win/tie/loss & unsafe \\
\midrule
small  & $6.2\%$  & $0.0\%$ & $92.4\%$ & $92.4\%$ & $\mathbf{93.5\%}$ & $+1.10$ & $0.62$ & 4/1/0 & 0 \\
medium & $54.5\%$ & $0.0\%$ & $40.0\%$ & $39.3\%$ & $\mathbf{45.5\%}$ & $+6.21$ & $0.00$ & 5/0/0 & 0 \\
large  & $72.4\%$ & $0.0\%$ & $15.2\%$ & $14.5\%$ & $\mathbf{25.0\%}$ & $+10.48$ & $5.86$ & 4/1/0 & 0 \\
\midrule
\multicolumn{6}{l}{\textbf{aggregate} (15 runs)} & $\mathbf{+5.93}$ & $5.07$ & \textbf{13/2/0} & \textbf{0} \\
\bottomrule
\end{tabular}
\caption{Work preserved, 56-agent testbed, five seeds per regime. $\Delta$ is
the \emph{paired} per-run difference in percentage points. \sys{} never lost
to \btwo{} in any run.}
\label{tab:main}
\end{table*}

Table~\ref{tab:main} is the main result. \sys{} preserved more correct work
than \btwo{} in 13 of 15 runs, tied in 2, lost in 0. The advantage grows with
the contaminated fraction, which is the direction the method predicts: the
more a reachability closure over-discards, the more there is to recover.

The two ties are informative and not noise. In both, selective recovery failed
verification and escalated to agent restart, at which point \sys{}'s discard
set becomes \btwo{}'s by construction. The ladder gives up the gain rather
than ship an uncertified recovery.

\subsection{Safety}

\begin{table}[t]
\centering
\small
\begin{tabular}{lrrr}
\toprule
 & \multicolumn{3}{c}{regime} \\
\cmidrule(lr){2-4}
quantity & small & medium & large \\
\midrule
unsafe preservations (all methods) & 0 & 0 & 0 \\
contamination escapes from closure & 0 & 0 & 0 \\
payload landed                     & 5/5 & 5/5 & 5/5 \\
recall                             & 1.0 & 1.0 & 1.0 \\
final task correct                 & 5/5 & 5/5 & 5/5 \\
degraded agents                    & 0 & 0 & 0 \\
\bottomrule
\end{tabular}
\caption{Safety, 15 runs. Every attack landed, so no zero here is arithmetic.}
\label{tab:safety}
\end{table}

Table~\ref{tab:safety} is the result we would retract the paper over. Across
all 15 runs and all four methods there were zero unsafe preservations.
Contamination never escaped the structural closure---an invariant the
comparison against \btwo{} depends on, tested here for the first time on a
topology with a broadcast hub and a memory channel that has no call-graph edge
back to its writer.

On the scripted matrix (96 cells, 30 repetitions, four detectors including a
blind control that flags nothing, and an adversarial self-reporter that lies
about which source moved its output), \sys{} recorded unsafe preservations in
\textbf{0 of its 24 cells}.

\subsection{Scripted matrix}

\begin{table}[t]
\centering
\small
\begin{tabular}{llrrr}
\toprule
scenario & variant & \btwo{} & \bone{} & \sys{} \\
\midrule
A web       & influencing   & $15.8\%$ & $21.1\%$ & $\mathbf{38.2\%}$ \\
A web       & exposed-only  & $15.8\%$ & $21.1\%$ & $\mathbf{90.4\%}$ \\
B memory    & influencing   & $57.9\%$ & $57.9\%$ & $57.9\%$ \\
B memory    & exposed-only  & $57.9\%$ & $57.9\%$ & $\mathbf{90.4\%}$ \\
C message   & influencing   & $15.8\%$ & $52.6\%$ & $\mathbf{57.9\%}$ \\
C message   & exposed-only  & $15.8\%$ & $52.6\%$ & $\mathbf{91.2\%}$ \\
\bottomrule
\end{tabular}
\caption{Work preserved, scripted mode, oracle detector, 30 repetitions.
Exposed-only variants plant a flagged source that instructs nothing---the case
where reachability baselines discard everything for nothing.}
\label{tab:scripted}
\end{table}

Table~\ref{tab:scripted} separates the two situations an operator faces. On
\emph{influencing} runs the payload really did change outputs and there is
genuinely less to save; \sys{} gains between $0$ and $22$ points. On
\emph{exposed-only} runs a flagged source was present but unused, and
reachability baselines discard $42$--$75$ points of work for nothing.

We draw attention to the one cell where \sys{} gains nothing: scenario B
influencing, where all three methods preserve $57.9\%$. We report it because
an earlier version of this work reported $63.2\%$ there, and the difference
was a defect (Section~\ref{sec:negative}).

\section{Economics: When This Is Worth It}
\label{sec:econ}

We now state the cost plainly, because the honest answer is not flattering and
a reader should not have to derive it.

\begin{table}[t]
\centering
\small
\begin{tabular}{lrrrr}
\toprule
regime & \btwo{} & \sys{} analysis & \sys{} total & vs \bzero{} \\
\midrule
small  & $1{,}238$ & $16{,}291$ & $17{,}374$ & $1.62\times$ \\
medium & $7{,}755$ & $15{,}766$ & $23{,}605$ & $2.21\times$ \\
large  & $9{,}272$ & $15{,}810$ & $25{,}082$ & $2.43\times$ \\
\bottomrule
\end{tabular}
\caption{Recovery cost in tokens, mean per run. \bzero{} (full restart) costs
$10{,}328$--$10{,}755$. \sys{} is more expensive than every baseline in every
regime.}
\label{tab:cost}
\end{table}

\subsection{\sys{} is not a token optimization}

Table~\ref{tab:cost} is unambiguous. \sys{} costs $2.7\times$--$14.0\times$
\btwo{} and $1.6\times$--$2.4\times$ a full restart. \textbf{If the only thing
that matters is tokens, and the workflow can simply be re-run, then restart is
the correct choice and our method is not.} We say so because any paper that
buries this invites the reviewer to find it.

The dominant term is attribution, not replay: analysis is $63\%$--$94\%$ of
\sys{}'s cost. Prior measurement on this system found that $68\%$ of
self-report calls land on events outside the contaminated closure and provably
cannot change the outcome; they are not skippable in the current design
because identifying them requires deferring attribution until after detection,
and deferral changes what the trace records. The overhead is therefore known
to be reducible in principle and unreduced in practice, which bounds how much
of Table~\ref{tab:cost} is intrinsic.

\subsection{When it is nonetheless the right choice}

The comparison in Table~\ref{tab:cost} assumes the preserved work is worth
exactly what it cost in tokens to produce. That assumption fails in three
common situations, and in those situations restart is not the cheap option but
the wrong one.

\paragraph{1. Work with side effects cannot be re-executed.}
An agent that sent a message, wrote to an external store, placed an order or
called a paid API has changed the world. ``Restart'' does not undo that; it
does it again. Preserving the clean $25\%$--$94\%$ of such a trace is not an
optimization, it is the only correct action. Our measurements do not quantify
this---our testbed's work is replayable by construction, which is what makes
it measurable---but it is the motivating deployment and the condition is
checkable statically.

\paragraph{2. Re-execution is not idempotent.}
LLM outputs are non-deterministic. A restart does not reproduce the original
run; it produces a different one that may be worse. In our frozen validation a
full restart \emph{failed the task} in one run where selective recovery
succeeded. Restart guarantees safety, not quality.

\paragraph{3. Latency and rate limits, not tokens, are the binding resource.}
In our runs, recovery wall-clock was dominated by external rate limiting, not
by token count. A method that re-executes fewer agents finishes sooner even
when it spends more tokens.

\subsection{A decision rule}

Let $N$ be the cost of re-executing the workflow, $A+R$ the cost of analysis
plus selective replay, $p$ the fraction of work preserved, and $V$ the value
per unit of preserved work. \sys{} is preferable to restart when
\[
p \cdot V \;>\; (A + R) - N .
\]
With our measured figures, the large regime requires $p=25.0\%$ of the trace
to be worth more than $\approx 14{,}750$ tokens, i.e.\ roughly $1.35$
workflow executions. Purely in tokens that is never satisfied. It is
satisfied whenever the preserved work is worth more than the tokens that
produced it---which is exactly the case when it has side effects, external
cost, or a quality that a re-run may not reproduce.

We consider stating this condition, rather than asserting a universal win, to
be part of the contribution.

\section{Negative Results}
\label{sec:negative}

Two results shaped this work and we report them because they are load-bearing.

\subsection{An ex-ante economic gate, falsified}

We built a gate to decide \emph{before} investigating whether investigation
could pay, of the form $\hat{A}/N + f_{\text{struct}} < 1$. On the 56-agent
testbed it was wrong in all three regimes, in both directions: it advised
INVESTIGATE where investigation bought nothing (spending $16{,}676$ tokens to
reach the same outcome a free restart reaches), and RESTART twice where
selective recovery preserved work a restart would have destroyed---including
one case where the restart failed the task and the selective recovery passed
it. Its estimator was off by $6.7\times$ on one regime.

The cause is that the scaling constant relating estimated to actual analysis
cost is a property of the client and prompt structure, not of the algorithm;
we measured values of $0.431$, $0.547$ and $1.781$ on three workloads of the
same system. We abandoned the gate rather than tune it, and the paper does not
depend on it.

\subsection{A safety defect that looked like a win}

Our contamination walk traversed influence and derivation edges. It did not
traverse \emph{ingestion}: the event that retrieves a source stores that
source's text in its own output, but the source does not exist when the event
is logged, so it is never in that event's exposures, so no verdict was ever
written for the pair and the walk never considered it. The event sat outside
the recovery region \emph{holding the payload}.

The verdict was not wrong; it was absent---which is worse, because every
consumer looks for a wrong verdict and none looks for a missing one.

Before the fix, the large regime read \sys{} $31.7\%$ preserved against
\btwo{}'s $14.5\%$: a $+17.2$ point win, with \textbf{5 unsafe
preservations}. After the fix it read $14.5\%$ with $0$. \emph{The entire
margin was the defect.} Had we stopped at that run, we would have published a
safety failure as a performance gain.

The general lesson is that preserved work and unsafe preservation are not
independent quantities: any mechanism that preserves more is mechanically a
candidate for preserving something it should not. The two must be read
together, and a gain accompanied by a non-zero unsafe count is not a gain. The
fix cost us $4.6$--$5.3$ points across ten scripted cells and one headline
claim, and we report the corrected numbers throughout.

\section{Limitations}

\begin{enumerate}
\itemsep2pt
\item \textbf{Three design points, replicated.} The five seeds per regime vary
model sampling, not the attack placement: the contaminated fraction is
identical across seeds within a regime, and the medium regime's paired
difference has standard deviation $0.00$. The effective sample for the
structural comparison is one workload per regime, confirmed five times. We
have since implemented seed-dependent placement and verified it produces five
distinct workloads per regime, but the campaign reported here is the fixed one.
\item \textbf{One topology.} Nothing here says how the gain behaves as hub
fan-out, stage depth or provider mix change.
\item \textbf{Oracle detector} in the real-system experiments. We measure
recovery after detection and say nothing about detection.
\item \textbf{The task is a copy task}, chosen so that capability failures of
a small local model do not masquerade as recovery failures. Nothing here shows
recovery works on hard tasks.
\item \textbf{No significance testing.} Five seeds per regime supports
descriptive statistics and a sign count, not a $p$-value.
\item \textbf{Cost is measured in tokens}, not money or latency, and
attribution is routed to the same model that produced the output; a cheaper
attribution model would change the economics and is untested.
\end{enumerate}

\section{Related Work}

Work on indirect prompt injection establishes the threat and its
channels~\cite{greshake2023,liu2024}, and a substantial literature addresses
detection and defence at the prompt or model
level~\cite{perez2022,chen2024}. Our work is downstream of all of it: we
assume detection and ask what to discard.

Provenance and taint tracking have a long history in systems
security~\cite{newsome2005,king2003}, where the standard discipline is
conservative reachability---precisely the \btwo{} baseline. The distinguishing
difficulty in the LLM setting is that presence in a context window is not
usage, and usage is not observable from the execution trace alone; it must be
established by asking the model or by counterfactual re-execution. The
asymmetric treatment of positive and negative evidence in
Section~\ref{sec:attribution} is our response to that.

Checkpoint-and-replay is standard in fault-tolerant
computing~\cite{elnozahy2002}; what differs here is that replay is not
required to be deterministic. Recovery never depends on a model reproducing a
prior output---it re-invokes only invalidated events and reuses logged bytes
everywhere else.

\section{Conclusion}

Exposure is not influence, and in a multi-agent system with any fan-out the
difference is large enough to matter: $13.1$ points of a 145-event trace in
our heaviest regime. \sys{} exploits that difference to preserve more correct
work than a reachability closure---in 13 of 15 runs, never fewer, with zero
unsafe preservations and recall $1.0$.

It is not cheaper. It costs $1.6\times$--$2.4\times$ a full restart in tokens,
and where work is freely re-executable we recommend restarting. Its value is
in the case where completed work cannot simply be redone, and there the
question is not whether recovery is worth two workflow executions but whether
the alternative---doing the side effects again---is acceptable at all.

We also report that our own ex-ante economic gate was falsified, and that a
provenance defect once presented itself to us as a $+17.2$ point improvement
while in fact preserving five contaminated events. We think a recovery paper
that cannot show it looked for that failure has not earned the reader's trust
in its safety numbers.

\begin{thebibliography}{9}
\small
\bibitem{greshake2023}
K.~Greshake et al.
\newblock Not what you've signed up for: Compromising real-world LLM-integrated
applications with indirect prompt injection.
\newblock In {\em AISec}, 2023.

\bibitem{perez2022}
F.~Perez and I.~Ribeiro.
\newblock Ignore previous prompt: Attack techniques for language models.
\newblock In {\em NeurIPS ML Safety Workshop}, 2022.

\bibitem{liu2024}
Y.~Liu et al.
\newblock Formalizing and benchmarking prompt injection attacks and defenses.
\newblock In {\em USENIX Security}, 2024.

\bibitem{chen2024}
S.~Chen et al.
\newblock Defending against prompt injection with structured queries.
\newblock In {\em USENIX Security}, 2024.

\bibitem{newsome2005}
J.~Newsome and D.~Song.
\newblock Dynamic taint analysis for automatic detection, analysis, and
signature generation of exploits on commodity software.
\newblock In {\em NDSS}, 2005.

\bibitem{king2003}
S.~T. King and P.~M. Chen.
\newblock Backtracking intrusions.
\newblock In {\em SOSP}, 2003.

\bibitem{elnozahy2002}
E.~N. Elnozahy, L.~Alvisi, Y.-M. Wang, and D.~B. Johnson.
\newblock A survey of rollback-recovery protocols in message-passing systems.
\newblock {\em ACM Computing Surveys}, 34(3):375--408, 2002.
\end{thebibliography}

\end{document}
